% 04-discussion.tex — Discussion, Limitations, and Conclusions
% Rule: Align claim strength with evidence (24 patients, LME non-convergence; avoid "prove", "demonstrate efficacy", "confirm"). Discuss all 6 limitations based on library/literature-matrix.md.

\section{Discussion}

\subsection{Principal findings}
Our longitudinal reanalysis of 356 fecal samples from 24 allo-HCT recipients illustrates how analytical model choice shapes the interpretation of clinical microbiome interventions. Across the cohort, transplantation was accompanied by a pronounced reduction in observed ASV richness between pre- and post-HCT time points ($H = 70.58$, $q = 1.33 \times 10^{-16}$), alongside substantial variation across stem cell graft sources and stool consistency categories. When all 356 samples were pooled across time in a cross-sectional Kruskal-Wallis test, no statistically significant difference in ASV richness was detected between the control ($n = 150$) and auto-FMT ($n = 206$) arms ($H = 1.29$, $p = 0.256$). In contrast, accounting for repeated within-patient measurements in a longitudinal linear mixed-effects (LME) model indicated a steep downward trajectory in the control group ($-0.859$ ASVs/day, $p < 0.001$) and a positive group-by-time interaction in the auto-FMT group ($+0.560$ ASVs/day, $p = 0.016$; implied net slope $-0.299$ ASVs/day), suggesting an attenuated rate of richness decline over the observation window. Concurrently, supervised random forest regression moderately predicted post-transplant week ($R^2 = 0.319$, $r = 0.565$, $p = 2.37 \times 10^{-7}$), identifying members of \textit{Enterococcaceae} (21.30\% combined importance across two features), \textit{Clostridium} (10.92\%), and \textit{Lachnospiraceae} (15.68\% combined across three features) as the primary taxa tracking temporal progression after allo-HCT.

The direction of the longitudinal interaction is consistent with the original trial, which reported that auto-FMT boosted microbial diversity \citep{taur2018reconstitution}. However, the original analysis included 25 patients, whereas our filtered reanalysis comprised 24 patients and 356 samples, and our single-slope LME on one richness index cannot by itself assess the compositional reestablishment described in that study.

\subsection{Comparison with previous methodological and clinical work}
The divergence between our null cross-sectional group comparison and the positive longitudinal interaction underscores a central principle in longitudinal microbiome research: pooling repeated observations across time conflates large within-subject temporal shifts with between-group treatment effects \citep{herman2024assessingb,zhang2026bayesian}. Because specimens collected from the same patient share host-specific baseline diversity and correlated temporal trajectories, treating all 356 samples as independent observations violates the exchangeability assumption of rank-based tests and allows the profound pre-to-post-HCT richness drop to obscure subtler differences in post-intervention recovery slopes \citep{herman2024assessingb,zhou2021linda}. In beta diversity ordinations, abundance-weighted phylogenetic distance (weighted UniFrac; PC1--PC3 cumulative variance 54.64\%) captured nearly twice the variation of presence/absence phylogenetic distance (unweighted UniFrac; 31.10\%) or non-phylogenetic Bray-Curtis dissimilarity (30.19\%), consistent with large shifts in dominant phylogenetic lineages---such as the post-transplant expansion of \textit{Enterococcaceae} and fluctuation of obligate anaerobic \textit{Clostridia} (\textit{Lachnospiraceae} and \textit{Clostridium})---driving community reorganization.

At the same time, recent methodological frameworks for fecal microbiota transplantation emphasize that neither alpha diversity recovery nor temporal shifts in taxonomic relative abundance are sufficient to establish donor microbiome engraftment \citep{herman2024assessing,herman2024assessingb}. As formalized by \citet{herman2024assessingb} and implemented in the QIIME 2 \texttt{q2-fmt} plugin \citep{herman2024assessing}, rigorous evaluation of FMT engraftment requires three complementary criteria: (1) chimeric asymmetric community coalescence, quantifying whether post-FMT recipient communities converge toward their specific baseline or donor reference profile over time; (2) donated microbiome indicator features, tracking the proportion of transferred donor features (such as Proportion of Engrafted Donated Features, PEDF) and differentially abundant taxa; and (3) temporal stability of the reconstituted community. Because our pipeline evaluated only cohort-level alpha diversity slopes and time-predictive genera without pairing post-FMT recipient profiles against each patient's own cryopreserved pre-antibiotic baseline, our findings reflect longitudinal richness retention and temporal succession rather than direct verification of autologous strain engraftment.

\subsection{Methodological limitations and future directions}
This reanalysis has six specific methodological limitations that temper clinical inference and motivate the adoption of specialized longitudinal microbiome tools:

\begin{enumerate}
  \item \textbf{Small sample size and unbalanced longitudinal design:} Although the analytical dataset comprised 356 fecal specimens, these observations were clustered within only 24 patients (10 control and 14 auto-FMT recipients), with per-patient sampling frequency ranging widely from 2 to 34 specimens (mean 14.8). This modest cohort size and irregular follow-up schedule limited statistical power and produced a wide 95\% confidence interval for the LME group-by-time interaction term ($0.104$ to $1.015$ ASVs/day, a nearly tenfold span). Future reanalyses should leverage estimation frameworks specifically designed for irregular, unbalanced longitudinal microbiome cohorts with missing visits \citep{barrera2025stochastic,fontaine2026locally}.
  \item \textbf{Non-convergence of the LME optimization algorithm:} The REML optimizer in \texttt{q2-longitudinal} reported non-convergence (\texttt{Converged: No}; residual scale $= 1,757.408$, patient random intercept variance $= 197.005$). In complex random-effects models for longitudinal microbiome data, estimation methods that rely on log-likelihood approximations are prone to biased estimates or unstable convergence \citep{barrera2025stochastic}. Consequently, the nominal $p = 0.016$ interaction effect must be interpreted cautiously as an exploratory trend rather than a definitive statistical conclusion. Applying a Stochastic Approximation Expectation-Maximization (SAEM) algorithm offers a more stable maximum likelihood estimation approach that directly accommodates unbalanced longitudinal designs without prior interpolation \citep{barrera2025stochastic}.
  \item \textbf{Oversimplification of pre- and post-FMT dynamics by a single linear slope:} Our LME specification fit a single constant linear slope across \texttt{day-relative-to-fmt} spanning days $-10$ to $+70$ relative to allo-HCT. Biologically, however, gut microbial richness undergoes a biphasic trajectory---rapid depletion during pre-FMT conditioning and antibiotic exposure followed by stabilization or recovery after auto-FMT reinfusion. Fitting one straight line across both phases averages pre-intervention injury with post-intervention recovery; thus, the net slope in the auto-FMT arm remained slightly negative ($-0.299$ ASVs/day) despite a positive interaction term ($+0.560$ ASVs/day). Piecewise (broken-stick) mixed models or locally sparse varying coefficient mixed models (LSVCMM), which exploit penalized kernel smoothing and random effects for serial correlation to identify specific time intervals of significant group differences \citep{fontaine2026locally}, would better match the underlying clinical timeline.
  \item \textbf{Sample loss and sampling variance from single-step rarefying:} Enforcing a uniform rarefying threshold of 10,000 reads per sample (\texttt{core-metrics-phylogenetic}) excluded 50 of the 406 in-window samples (12.32\%) and discarded excess sequencing reads from deeper libraries. As demonstrated by \citet{mcmurdie2013waste}, single rarefying is statistically inadmissible because it discards valid information and inflates Type-II error rates; moreover, in allo-HCT cohorts, low-depth specimens may disproportionately represent severe post-antibiotic bacterial depletion, introducing potential selection bias. Within the QIIME 2 ecosystem, replacing single rarefying with bootstrapped or multi-rarefaction diversity averaging via \texttt{q2-boots} \citep{raspet2024facilitating} would incorporate all valid sequence data and quantify subsampling uncertainty.
  \item \textbf{Pseudoreplication in cross-sectional rank tests:} The pairwise Kruskal-Wallis comparisons treated 356 correlated observations from 24 patients as independent samples. Ignoring within-subject correlation in repeated-measures microbiome data is known to distort type I and type II error rates \citep{herman2024assessingb,zhou2021linda,zhang2026bayesian}. Cross-sectional rank tests should therefore be restricted to single-time-point summaries (e.g., baseline or a fixed landmark week) or replaced by mixed-effects models that explicitly partition within- and between-subject variance.
  \item \textbf{Temporal prediction versus therapeutic engraftment and compositional bias:} In our feature volatility analysis, random forest regression was trained on genus-level relative frequencies to predict \texttt{week-relative-to-hct} ($R^2 = 0.319$), without stratifying by treatment arm or correcting for compositional bias. Consequently, top-ranked taxa such as \textit{Enterococcaceae} and \textit{Lachnospiraceae} represent markers of time since transplantation rather than biomarkers of auto-FMT efficacy. Furthermore, because amplicon sequencing yields compositional relative abundances with high zero inflation, future analyses should apply centered log-ratio (CLR) linear mixed models with mode-based bias correction (\texttt{LinDA-LMM}) \citep{zhou2021linda} or Bayesian nonparametric longitudinal factor models (\texttt{LODI}) \citep{zhang2026bayesian} to test group-by-time differential abundance while accounting for dynamic microbial co-abundance networks.
\end{enumerate}

\subsection{Conclusions}
In this longitudinal QIIME 2 reanalysis of 24 allo-HCT patients (356 fecal samples), gut microbial richness declined sharply after transplantation, and supervised random forest regression identified \textit{Enterococcaceae}, \textit{Clostridium}, and \textit{Lachnospiraceae} as key indicators of post-transplant temporal succession ($R^2 = 0.319$). While pooled cross-sectional Kruskal-Wallis testing showed no difference between randomized arms ($p = 0.256$), longitudinal LME modeling detected a slower rate of richness decline in the auto-FMT group ($+0.560$ ASVs/day interaction, $p = 0.016$). However, because the cohort was small and unbalanced, the LME optimizer did not converge, a single linear slope was fit across pre- and post-FMT windows, and donor strain coalescence was not directly evaluated, these results should be viewed as methodological and exploratory rather than conclusive evidence of clinical efficacy. Incorporating multi-rarefaction diversity (\texttt{q2-boots}), patient-baseline coalescence metrics (\texttt{q2-fmt}), stochastic mixed-model estimation (SAEM), time-varying coefficient models (LSVCMM), and compositionally adjusted longitudinal differential abundance (\texttt{LinDA-LMM}) will be essential for robust inference in future microbiome intervention trials.
